% ================================================================
% SECTION 18: PROBLEM IMPLEMENTATIONS
% ================================================================
\section{Problem Implementations}
\label{sec:problems}

\begin{center}
\code{bionetflux.problems.ks\_problem} \quad(569 lines) \\
\code{bionetflux.problems.ooc\_problem} \quad(596 lines) \\
\code{bionetflux.utils.mesh\_mapping} \quad(100 lines)
\end{center}

Each problem module is a \emph{recipe}: a standalone Python file that,
given an optional geometry and an optional TOML file, produces the
four objects expected by \code{SolverSetup.initialize()}.

% ------------------------------------------------------------------
\subsection{Module-Level Entry Point}

Both modules expose:

\begin{lstlisting}[language=Python, caption={Shared module entry point}]
def create_global_framework(
    geometry: Optional[DomainGeometry] = None,
    config_file: Optional[str] = None
) -> Tuple[List[Problem], GlobalDiscretization,
           ConstraintManager, str]:
\end{lstlisting}

If \code{geometry} is \code{None}, a built-in default geometry is created.  If
\code{config\_file} is \code{None}, the config manager's defaults apply.

\subsubsection{Internal Workflow (common to both)}

\begin{enumerate}
  \item Load and validate the TOML configuration via the problem-specific config manager
\\    (\code{KSConfigManager} or \code{OoCConfigManager}).
  \item Build or accept the \code{DomainGeometry}.
  \item For each domain in the geometry, create:
    \begin{itemize}
      \item A \code{Problem} instance with physical parameters, initial conditions, forcing functions, and (optionally) exact solutions as callables.
      \item A \code{Discretization} from \code{n\_elements} (or from target $h$ via \code{compute\_n\_elements\_from\_h}).
    \end{itemize}
  \item Assemble a \code{GlobalDiscretization} and call \code{set\_time\_parameters(dt,\,T)}.
  \item Build a \code{ConstraintManager}:
    \begin{itemize}
      \item Homogeneous Neumann BCs at all exterior boundary points.
      \item Trace-continuity (or Kedem-Katchalsky) at all interior junctions.
    \end{itemize}
  \item Optionally apply per-point boundary overrides from the TOML \code{[boundary\_conditions]}
    section via \code{apply\_boundary\_overrides}.
  \item Call \code{constraint\_manager.map\_to\_discretizations()}.
  \item Return \code{(problems, global\_discretization, constraint\_manager, problem\_name)}.
\end{enumerate}

% ------------------------------------------------------------------
\subsection{Keller-Segel Module}

\begin{longtable}{p{5cm}p{9.5cm}}
\toprule
\textbf{Key} & \textbf{Value} \\
\midrule
Config type & \code{"ks"} \\
Core type & \code{"keller\_segel"} (for \code{StaticCondensationFactory}) \\
$n_\mathrm{eq}$ & 2: $u$ (cell density), $\varphi$ (chemical signal) \\
Default geometry & \code{build\_arc\_sequence\_geometry(N=1, start=2.0, length=1.0)} \\
Parameters & $\mu$, $\nu$, $a$, $b$ \quad plus callables $\chi(s,t)$, $\chi'(s,t)$ \\
\bottomrule
\end{longtable}

Additional helper functions:
\begin{itemize}
  \item \code{build\_default\_ks\_geometry()} --- one-arc default.
  \item \code{setup\_constraints\_from\_geometry(geometry, problems, neq)} --- constraint builder used internally.
\end{itemize}

% ------------------------------------------------------------------
\subsection{Organ-on-Chip Module}

\begin{longtable}{p{5cm}p{9.5cm}}
\toprule
\textbf{Key} & \textbf{Value} \\
\midrule
Config type & \code{"ooc"} \\
Core type & \code{"organ\_on\_chip"} (for \code{StaticCondensationFactory}) \\
$n_\mathrm{eq}$ & 4: $u$, $\omega$, $v$, $\varphi$ \\
Default geometry & \code{build\_grid\_geometry()} \\
Parameters & $\nu$, $\mu$, $\epsilon$, $\sigma$, $a$, $b$, $c$, $d$, $\chi$  plus callables \\
\bottomrule
\end{longtable}

Additional helpers:
\begin{itemize}
  \item \code{build\_default\_geometry()} --- default OoC grid.
  \item \code{build\_T\_junction\_geometry()} --- alternative T-junction layout.
\end{itemize}

% ------------------------------------------------------------------
\subsection{Writing a Custom Problem Module}

Any Python file that exposes \code{create\_global\_framework(geometry, config\_file)}
with the same return signature can serve as a problem module.  The
minimum steps are:

\begin{enumerate}
  \item Define or subclass a \code{BaseConfigManager} for TOML support.
  \item Register the corresponding \code{StaticCondensation} class via \\
        \code{StaticCondensationFactory.register\_implementation("my\_type", MyStaticCond)}.
  \item Build \code{Problem} objects with appropriate \code{type="my\_type"}.
  \item Return the standard 4-tuple.
\end{enumerate}

% ------------------------------------------------------------------
\subsection{Mesh Mapping Utilities}
\label{sec:mesh_mapping}

\begin{center}
\code{bionetflux.utils.mesh\_mapping}
\end{center}

\begin{lstlisting}[language=Python, caption={Mesh mapping utilities}]
def parametric_to_physical_mesh(domain: DomainInfo,
                                discretization: Discretization
                                ) -> Tuple[np.ndarray, np.ndarray]:
    """Map parametric nodes to (x, y) physical coordinates."""

def create_physical_mesh_dict(domain, discretization) -> dict:
    """Full mesh dictionary with nodes_physical, nodes_parametric, etc."""
\end{lstlisting}

Both functions perform a linear interpolation between the parametric coordinates of the extrema \code{domain.extrema\_start}
and \code{domain.extrema\_end} using the normalised parametric coordinate
$t = (s - s_0) / L$.
